\documentclass[10pt,a4paper]{article}
\usepackage[utf8]{inputenc}
\usepackage[T1]{fontenc}
\usepackage[margin=18mm]{geometry}
\usepackage{longtable}
\usepackage{hyperref}
\usepackage{enumitem}
\setlist[itemize]{noitemsep, topsep=2pt}
\setlength{\parindent}{0pt}
\begin{document}

\section*{Non-Verbal \& Multimodal Discourse Analysis Worksheet}

\textit{Qualitative analysis of pauses, trailing sentences, prosody/intonation, vocalisations (coughs, sighs, scoffs), and kinesics (hand gestures, facial expressions, posture) in political news discourse}\par\smallskip

\bigskip\hrule\bigskip

\subsection*{Guidelines: Five non-verbal \& paralinguistic features to look for in every event}

\begin{longtable}{|p{0.23\textwidth}|p{0.23\textwidth}|p{0.23\textwidth}|p{0.23\textwidth}|}
\hline
Feature & What to look for & Illustrative example (from sample) & Questions to answer \\ \hline \endhead
\textbf{1. Pauses, Silences \& Abrupt Stops} & Micro-pauses \texttt{(.)}, timed silences \texttt{(1.5)}, pregnant/dramatic pauses before a loaded word, hesitant stumbles, abrupt self-interruptions or cut-offs (\texttt{---}), and forced silences (muting a panelist's mic or talking over them). & Anchor pauses \texttt{(1.5)} before saying \textit{"'loota'"} for dramatic weight; Addl. DCP cuts off mid-sentence (\textit{"wo toh--- abhi investigation chal raha hai"}) when asked if police struck first; panelist's mic faded during counter-question. & Where does the pause or stop occur? Is it deliberate/rhetorical (for suspense, emphasis, or gravitas) or defensive/evasive (cognitive load, self-repair, avoiding blame)? Who controls turn-taking and silence? \\ \hline
\textbf{2. Trailing Sentences \& Unfinished Utterances} & Clauses that trail off (\texttt{...}), drop in volume at the end, or leave the conclusion unspoken (aposiopesis) so the audience supplies the implication without the speaker taking explicit responsibility. & \textit{"Aur phir wahan jo log ghuse the... aap samajh hi sakte hain..."}; \textit{"500 salon ke sangharsh ke baad agar aisa..."} trailing off with a slow head shake. & What exact thought is left unfinished? What presupposition or shared bias is the viewer invited to fill in? Does trailing off act as a safe hedge to imply higher-level complicity or "anti-national" motive without stating it? \\ \hline
\textbf{3. Intonation, Pitch, Volume \& Prosody} & Pitch contour (rising \texttt{$\uparrow$} for incredulity/challenge; falling \texttt{$\downarrow$} for finality), volume shifts (shouting vs. hushed whisper), emphatic word stress, speech rate (rapid-fire urgency vs. slow solemnity), and tonal register (sarcastic, reverential, combative). & Hushed, reverential cadence on \textit{"Prabhu Shri Ram ke Mandir"} vs. sharp, high-pitched interrogation on \textit{"par Akhilesh ji Ram Mandir kyon$\uparrow$ nahin jate?!"}; satirical sing-song inflection on \textit{"priventiv arenjment"}. & Which exact words receive pitch peaks or volume stress? How does vocal tone turn neutral words into irony, mockery, or moral outrage? Does the speaker shift prosodic register depending on which actor is discussed? \\ \hline
\textbf{4. Coughs, Sighs \& Non-Lexical Vocalisations} & Audible sighs (\texttt{[sigh]}), scoffs/snorts (\texttt{[scoff]}), ironic chuckles (\texttt{[chuckle]}), throat-clearing/coughs (\texttt{[cough]}), sharp intakes of breath/gasps (\texttt{[gasp]}), and tongue-clicks (\texttt{[tsk]}). & Heavy \texttt{[sigh]} before reporting Trust oversight lapses (signalling reluctant sorrow rather than hostility); mocking \texttt{[scoff]} or \texttt{[chuckle]} when reading an opponent's denial; \texttt{[cough / throat-clear]} when pivoting away from an awkward fact. & What emotion or epistemic stance does the vocalisation perform (exasperation, moral disgust, disbelief, ridicule, discomfort)? Does a sigh or scoff pre-judge a speaker's claim before any verbal rebuttal is made? \\ \hline
\textbf{5. Hand Gestures, Facial Expressions \& Posture (Kinesics)} & Hand movements (air-quotes, finger-pointing/wagging, desk-thumping, open palms, counting on fingers, dismissive wave), facial cues (raised eyebrows, eye-rolls, smirk, furrowed brow), head nods/shakes, leaning into camera, and visual props (bandaged limbs, documents). & Anchor forms physical air-quotes on \textit{"'loota'"} or \textit{"'peaceful' protest"}; host points finger/pen at lens during an ultimatum; injured DCP holds bandaged hand up to camera; creator raises one eyebrow and leans in for a conspiratorial reveal. & How do gestures reinforce, contradict, or sharpen the spoken words? During split-screen listening, does the host's face (nodding vs. smirking/eye-rolling) cue the audience on whom to trust? How are bodily visuals used as evidence? \\ \hline
\end{longtable}

\subsection*{Quick Multimodal Transcription Notation Key (use inside transcript excerpts)}

\begin{longtable}{|p{0.15\textwidth}|p{0.15\textwidth}|p{0.15\textwidth}|p{0.15\textwidth}|p{0.15\textwidth}|p{0.15\textwidth}|}
\hline
Symbol & Meaning & Example in transcript & Symbol & Meaning & Example in transcript \\ \hline \endhead
\texttt{(.)} & Brief micro-pause (< 0.5 sec) & \texttt{woh (.) suparavijan lacking tha} & \texttt{$\uparrow$ / $\downarrow$} & Marked rising / falling pitch & \texttt{kyon$\uparrow$ FIR nahin karaee$\downarrow$} \\ \hline
\texttt{(1.5)} & Timed pause in seconds & \texttt{kisne (1.5) 'loota'?} & \texttt{WORD / \_word\_} & Loud volume / emphatic stress & \texttt{ye PAAP hai / \_bilkul\_ galat} \\ \hline
\texttt{---} & Abrupt stop, cut-off, or self-repair & \texttt{police ki taraf se--- abhi jaanch...} & \texttt{$^\circ$word$^\circ$} & Softened / whispered delivery & \texttt{$^\circ$ander ki baat ye hai$^\circ$} \\ \hline
\texttt{...} & Trailing off / unfinished sentence & \texttt{aap jaante hain ki kaun log...} & \texttt{[sigh] [cough] [scoff]} & Non-lexical vocalisations & \texttt{[sigh] 500 salon ke baad...} \\ \hline
\texttt{[overlap]} & Simultaneous / overlapping speech & \texttt{[Anjana: suniye---] [Guest: nahin!]} & \texttt{\{\{gesture: ...\}\}} & Hand gesture / facial cue / posture & \texttt{\{\{air-quotes on 'loota'\}\}} \\ \hline
\end{longtable}

\textbf{Methodological Reminder:} A non-verbal cue is not automatically manipulative or ideological---a cough, pause, or rapid cadence may reflect studio audio delay, live time pressure, or a speaker's habitual style. Always check (1) whether the cue clusters with verbal features, and (2) whether the speaker's delivery shifts noticeably across different actors or events (within-outlet baseline).\par\smallskip

\bigskip\hrule\bigskip

\subsection*{Event 1. Ram Mandir donation controversy (Non-Verbal \& Paralinguistic Layer)}

\begin{itemize}
  \item \textbf{Focus:} How do prosody, pauses, trailing sentences, vocalisations (sighs, coughs, scoffs), and body language frame the controversy---as solemn devotional sorrow, administrative scrutiny, or partisan combat?
  \item \textbf{Outlet:} Mohak Mangal (\textit{Soch by Mohak Mangal}) | \textbf{Date:} N/A (Skipped per project guidelines)
  \item \textbf{Video title:} \textit{Skipped --- No video relevant to the Ram Mandir donation/land scam; channel only published a historical explainer on the Ayodhya dispute} | \textbf{Timestamp(s):} N/A
  \item \textbf{Analyst:} Baani Solanki (\texttt{2025114016}) | \textbf{Language:} N/A (Skipped)
  \item \textbf{Format / Visual setup:} \texttt{[ ]} TV Panel Debate \texttt{[ ]} TV Bulletin \texttt{[x]} Solo YouTube Explainer \texttt{[ ]} Interview / Discussion
  \item \textbf{Primary speaker(s) observed:} N/A (Skipped for Event 1 per official project guidelines: \textit{"Mohak mangal - no video relevant to scam; there is video regarding history, so skipping"})
\end{itemize}

\subsubsection*{Event-specific non-verbal questions (Event 1)}

\noindent\textbf{1.} \textbf{Vocal \& gestural marking of core labels ("scam", "loota", "paap", "gadabdi", "avyavastha"): When the controversy is named, what non-verbal cues accompany the label? (Look for physical air-quotes, vocal scare-quotes, dramatic pauses before the word, hushed/reverential drops in volume, or shouted stress.)}\par\smallskip
\begin{itemize}
  \item \textbf{Notes:}
\textbf{Skipped per project instructions.} From a discourse-analytic standpoint, Mohak Mangal's non-coverage of the 2021/2024 Ram Mandir donation and land-deal allegations constitutes a \textbf{structural silence (editorial gatekeeping)}: unlike Dhruv Rathee (who foregrounds day-to-day political/corruption exposés) or TV news channels (which run daily partisan debates), \textit{Soch by Mohak Mangal} avoids acute partisan flashpoints involving religious institutions unless they crystallise into broader policy or structural case studies.\par\smallskip
\end{itemize}

\noindent\textbf{2.} \textbf{Differential demeanour toward actors (Temple, Trust leadership, counting clerks, donors, opposition accusers): How do the speaker's voice pitch, facial expression, and posture shift when referring to Prabhu Shri Ram / devotees vs. top figures (Champat Rai, Nripendra Misra, CM Yogi) vs. arrested counting staff vs. opposition leaders (Akhilesh Yadav, Awadhesh Prasad)?}\par\smallskip
\begin{itemize}
  \item \textbf{Notes:}
\textbf{Skipped per project instructions} (no video covering Event 1).\par\smallskip
\end{itemize}

\noindent\textbf{3.} \textbf{Solemn devotional grief vs. high-decibel partisan spectacle: How do pauses, sighs, speech rate, and hand gestures construct the segment's emotional register? Where is delivery slow, hushed, and mournful (sighs, bowed head, folded hands), and where does it escalate into rapid-fire shouting, desk-thumping, or mocking irony?}\par\smallskip
\begin{itemize}
  \item \textbf{Notes:}
\textbf{Skipped per project instructions} (no video covering Event 1).\par\smallskip
\end{itemize}

\noindent\textbf{4.} \textbf{Non-verbal evaluation of financial scrutiny: When questions about unreceipted offerings, lack of frisking, or land deals are raised, how does the host/speaker react non-verbally (earnest nodding, leaning into the camera, open-palm appeals vs. eye-rolls, scoffs, dismissive hand waves, or incredulous rising pitch $\uparrow$)?}\par\smallskip
\begin{itemize}
  \item \textbf{Notes:}
\textbf{Skipped per project instructions} (no video covering Event 1).\par\smallskip
\end{itemize}

\noindent\textbf{5.} \textbf{Performing certainty vs. hedging and innuendo (pauses, stops, trailing off): Are allegations delivered with definitive falling intonation $\downarrow$, finger-counting, and pointing at on-screen evidence (CCTV, documents), or hedged with hesitant pauses (.), abrupt self-interruptions (---), throat-clearing [cough], and trailing sentences (...) that imply higher-level complicity without naming it?}\par\smallskip
\begin{itemize}
  \item \textbf{Notes:}
\textbf{Skipped per project instructions} (no video covering Event 1).\par\smallskip
\end{itemize}

\noindent\textbf{6.} \textbf{Turn-taking, interruptions, and acoustic treatment of denials: In debates/interviews: who receives calm, uninterrupted speaking time and sympathetic nods, and who is cut off mid-sentence (---), talked over, or muted? In solo videos: how are clips of official denials framed via reaction pauses, raised eyebrows, or sarcastic chuckles?}\par\smallskip
\begin{itemize}
  \item \textbf{Notes:}
\textbf{Skipped per project instructions} (no video covering Event 1).\par\smallskip
\end{itemize}

\noindent\textbf{7.} \textbf{Performing the boundary between religious devotion and institutional critique: How does the presenter use vocal tone, sighs, and body language to signal personal reverence for the Temple (or secular objectivity) while criticising---or defending---the Trust's financial management?}\par\smallskip
\begin{itemize}
  \item \textbf{Notes:}
\textbf{Skipped per project instructions} (no video covering Event 1).\par\smallskip
\end{itemize}

\noindent\textbf{8.} \textbf{Self-presentation as the "guardian of public trust": How does the speaker use direct eye-gaze, posture, moral sighs, or emphatic gestures (pointing at the viewer, hands on heart/desk) to position themselves as the authentic protector of devotees or public accountability?}\par\smallskip
\begin{itemize}
  \item \textbf{Notes:}
\textbf{Skipped per project instructions} (no video covering Event 1).\par\smallskip
\end{itemize}

\subsubsection*{Apply the five non-verbal features (Event 1)}

\begin{longtable}{|p{0.31\textwidth}|p{0.31\textwidth}|p{0.31\textwidth}|}
\hline
Non-verbal feature & Timestamp + spoken anchor words + non-verbal notation & What it does in context (stance, framing, implication) \\ \hline \endhead
\textbf{1. Pauses, silences \& stops} & \texttt{N/A (Event 1 skipped per project guidelines)} & Editorial omission / structural silence on acute religious-trust financial allegations; channel focuses on long-horizon policy and sociological explainers (see Event 2). \\ \hline
\textbf{2. Trailing sentences (\texttt{...})} & \texttt{N/A (Event 1 skipped per project guidelines)} & N/A \\ \hline
\textbf{3. Intonation, pitch \& volume} & \texttt{N/A (Event 1 skipped per project guidelines)} & N/A \\ \hline
\textbf{4. Coughs, sighs \& vocalisations} & \texttt{N/A (Event 1 skipped per project guidelines)} & N/A \\ \hline
\textbf{5. Hand gestures, face \& posture} & \texttt{N/A (Event 1 skipped per project guidelines)} & N/A \\ \hline
\end{longtable}

\subsubsection*{Short summary (Event 1 Non-Verbal Synthesis)}

\begin{itemize}
  \item \textbf{What overall interpretation or emotional stance does the non-verbal delivery encourage?}
\textbf{Skipped per project guidelines.} Mohak Mangal did not produce a video on the Ram Mandir donation/land controversy.\par\smallskip
  \item \textbf{Whose voice or physical presence is acoustically/visually foregrounded, undermined, or silenced?}
\textbf{Skipped per project guidelines.}\par\smallskip
  \item \textbf{What evidence supports your non-verbal reading? What complicates it (e.g., speaker baseline, format)?}
\textbf{Skipped per project guidelines.} Comparing Mohak Mangal's absence in Event 1 with his deep-dive policy explainer in Event 2 confirms his channel's editorial baseline: avoiding breaking-news partisan allegations around religious bodies while investing heavily in educational, sociological, and institutional-reform topics affecting young Indians.\par\smallskip
\end{itemize}

\bigskip\hrule\bigskip

\subsection*{Event 2. Students' protest regarding NEET paper leak (Non-Verbal \& Paralinguistic Layer)}

\begin{itemize}
  \item \textbf{Focus:} How do prosody, pauses, trailing sentences, vocalisations, and visual body language construct the legitimacy of student protesters, police use of force, and state accountability?
  \item \textbf{Outlet:} Mohak Mangal (\textit{Soch by Mohak Mangal}) | \textbf{Date:} 14 June 2024
  \item \textbf{Video title:} NTA: The Biggest Scam in Indian Education? | NEET 2024 | Mohak Mangal | \href{https://youtu.be/v-RikSqb0Vw}{https://youtu.be/v-RikSqb0Vw} | \textbf{Timestamp(s):} \texttt{[00:00]}--\texttt{[20:12]}
  \item \textbf{Analyst:} Baani Solanki (\texttt{2025114016}) | \textbf{Language:} Hinglish (Hindi in Roman script + English terms)
  \item \textbf{Format / Visual setup:} \texttt{[ ]} TV Panel Debate \texttt{[ ]} TV Bulletin \texttt{[x]} Solo YouTube Explainer \texttt{[ ]} Interview / Discussion
  \item \textbf{Primary speaker(s) observed:} Mohak Mangal (solo studio presenter seated/standing in front of warm minimalist studio lighting, integrated with data visualisations, animated score graphs, and clips of NEET students and NTA press conferences)
\end{itemize}

\subsubsection*{Event-specific non-verbal questions (Event 2)}

\noindent\textbf{1.} \textbf{Prosodic \& gestural marking of participant labels ("students/bachche" vs. "protesters" vs. "upadravi/goons"): How do vocal tone, volume, facial expression, and hand gestures shift across labels? (e.g., softened, empathetic pitch for "students" vs. harsh, staccato stress, scowls, or accusatory finger-pointing on "upadraviyon", "goons", or "anti-social elements".)}\par\smallskip
\begin{itemize}
  \item \textbf{Notes:}
Mohak Mangal never uses derogatory labels (\textit{"upadravi"}, \textit{"goons"}). He refers to candidates by personal names---such as \texttt{"Ayushi"} (\texttt{[00:15]}) and \texttt{"Anjali"} (\texttt{[01:02]})---or as \textit{"students"}, \textit{"bachche"}, and \textit{"middle-class families"} (\texttt{[01:45]}, \texttt{[16:20]}). When introducing individual students' stories (\texttt{[00:00]}--\texttt{[02:15]}), his voice drops into an \textbf{intimate, narrative, low-decibel register} (\texttt{$^\circ$word$^\circ$}), his eyebrows draw together in quiet empathy, and his hands rest calmly or open gently toward the camera \texttt{\{\{gesture: gentle open-palm extension\}\}}. When referring to the \textit{"NTA"} or \textit{"coaching nexus"}, he does not shout; instead, his tone becomes \textbf{clinical, measured, and dryly analytical}, using precision thumb-and-index-finger pinch gestures \texttt{\{\{gesture: thumb-index precision pinch\}\}} to dissect institutional failure.\par\smallskip
\end{itemize}

\noindent\textbf{2.} \textbf{Delivery of students' stated demands: When the protesters' core demands (Education Minister's resignation, compensation, release of Sonam Wangchuk, exam reform) are mentioned, are they delivered slowly with deliberate pauses and finger-counting, rushed through in a flat monotone, mocked with air-quotes/scoffs, or trailed off (...)?}\par\smallskip
\begin{itemize}
  \item \textbf{Notes:}
Mohak delivers the students' demands and his own structural reform proposals (\texttt{[16:40]}--\texttt{[19:30]})---(1) independent external audit of NTA's grace-mark algorithm, (2) investigation of the Bihar/Godhra paper leak centres, and (3) decentralising/reforming India's single-point high-stakes exam architecture---\textbf{slowly, with deliberate \texttt{(0.8)} pauses and structured finger-counting} \texttt{\{\{gesture: counting 1, 2, 3 on fingers\}\}}. Each point concludes with a calm, definitive falling pitch \texttt{$\downarrow$} and an earnest nod, treating student demands not as street slogans but as serious public-policy imperatives.\par\smallskip
\end{itemize}

\noindent\textbf{3.} \textbf{Where does vocal \& kinetic energy peak---grievance, street disruption, or electoral horse-race?: Track where the speaker's volume, pitch range, pace, and gesture intensity reach their climax: during accounts of the paper leak/student suicides, during street clashes and barricade breaches, or during partisan BJP-vs-Opposition sparring.}\par\smallskip
\begin{itemize}
  \item \textbf{Notes:}
Because Mohak's baseline is exceptionally composed (\textasciitilde{}125--135 WPM, seminar-style delivery), his energy peaks through \textbf{controlled prosodic contrast and emotional gravity} rather than high-decibel shouting:\par\smallskip
  \item \textbf{Peak 1 (Statistical Proof of Rank Inflation, \texttt{[04:10]}--\texttt{[07:45]}):} His pace quickens and his pitch rises markedly \texttt{$\uparrow$} when demonstrating on a graph how a score of \texttt{620} dropped from rank \texttt{\textasciitilde{}15,000} in 2023 to rank \texttt{\textasciitilde{}60,000} in 2024 (\texttt{"yani \_chaar\_ guna rank gir gayi$\uparrow$!"}, \texttt{[05:12]}), accompanied by a sweeping downward hand trajectory \texttt{\{\{gesture: right hand sweeping sharply downward along chart\}\}}.
  \item \textbf{Peak 2 (Human Cost \& Trust Erosion, \texttt{[17:15]}--\texttt{[19:10]}):} The emotional climax occurs when he discusses the psychological collapse of \texttt{24 lakh} young Indians losing faith in meritocracy (\texttt{"jab mehnat ka fal nahin milta..."}, \texttt{[18:05]}), where he slows down to \textasciitilde{}105 WPM, pauses \texttt{(1.5)} seconds, and looks silently into the lens.
\end{itemize}

\noindent\textbf{4.} \textbf{Non-verbal attribution of responsibility (sighs, stops, head shakes, pointing): When examination authorities (NTA), the Education Minister, police, opposition politicians, or students are named as responsible, what non-verbal cues accompany the attribution (exasperated sighs, mocking laughs, abrupt self-stops ---, or direct finger-pointing at the lens)?}\par\smallskip
\begin{itemize}
  \item \textbf{Notes:}
Responsibility is attributed primarily to the \textbf{National Testing Agency's (NTA) structural opacity and incompetence} through understated, devastating paralinguistic cues:\par\smallskip
  \item When playing clip excerpts of NTA Director General Subodh Kumar Singh claiming at the press conference that \textit{"everything is transparent, integrity is not compromised"} (\texttt{[08:50]}), Mohak cuts back to himself, pauses \texttt{(1.2)}, lets out a quiet, weary \texttt{[sigh]}, and gives a slow head shake \texttt{\{\{gesture: slow bilateral head shake + pursed lips\}\}} before saying: \textit{"Yeh press conference dekh kar (.) answers kam mile (.) aur sawaal \_zyada\_ khade ho gaye$\downarrow$"} (\texttt{[09:12]}).
  \item When noting that the head of the NTA's "grievance committee" is the UPSC Chairman/same bureaucratic circle (\texttt{[09:40]}), Mohak raises both eyebrows and gives a brief, dry \texttt{[chuckle]} of disbelief.
\end{itemize}

\noindent\textbf{5.} \textbf{Prosody of questions---open inquiry vs. loaded rhetorical challenge: Are questions about the protest delivered with level, open intonation, or with sharp rising pitch $\uparrow$, sarcastic stress, raised eyebrows, and trailing clauses (...) that presuppose protesters are violent or politically manipulated?}\par\smallskip
\begin{itemize}
  \item \textbf{Notes:}
Mohak uses \textbf{Socratic, analytical inquiry} to guide the viewer step-by-step. When asking whether the anomaly was caused by an easier paper, grace marks, or a paper leak (\texttt{[03:30]}), he uses level, open-minded inquiry intonation with upturned palms: \textit{"Chaliye (.) teeno possibilities ko data ke saath check karte hain$\downarrow$"} (\texttt{[03:38]}). However, once the data eliminates NTA's excuses, his questions toward the NTA take on a \textbf{pointed, incredulous rising contour \texttt{$\uparrow$}}: \textit{"Agar grace marks sirf \texttt{1,563} students ko mile (.) toh baaki har score bracket mein \texttt{30,000} se \texttt{40,000} rank ka inflation kaise$\uparrow$ aaya?!"} (\texttt{[07:18]}).\par\smallskip
\end{itemize}

\noindent\textbf{6.} \textbf{Visual \& paralinguistic performance of injury and violence: How are claims of violence substantiated non-verbally on screen? (Compare camera close-ups and pained vocal tone on police bandages/fractured hands vs. visual/acoustic treatment of student head injuries, lathi-charge footage, or defensive stammers/pauses when an officer or guest is asked who used force first.)}\par\smallskip
\begin{itemize}
  \item \textbf{Notes:}
Mohak frames injury in \textbf{psychological, familial, and socio-economic terms} rather than street violence. He shows close-up B-roll of exhausted students sitting outside coaching centres in Kota and Delhi, parents mortgaging land for coaching fees, and peaceful candlelight/placard demonstrations (\texttt{[01:15]}, \texttt{[15:50]}). His vocal tone during these visuals is hushed and compassionate (\texttt{$^\circ$ek middle-class parivar apni poori jama-poonji laga deta hai$^\circ$}, \texttt{[01:32]}), accompanied by a solemn head nod.\par\smallskip
\end{itemize}

\noindent\textbf{7.} \textbf{Asymmetric delivery of police actions vs. protester actions: Compare the vocal cadence, pauses, and hand gestures used when describing police actions (e.g., calm, matter-of-fact administrative tone on "tear gas chhode gae" / "preventive arrangement") versus protester actions (alarmist pitch, striking gestures, or ironic softening as "retaliation").}\par\smallskip
\begin{itemize}
  \item \textbf{Notes:}
Mohak contrasts the \textbf{earnest, evidence-backed mobilisation of students and teachers} with the \textbf{evasive bureaucracy of the NTA}:\par\smallskip
  \item \textbf{Students and educators} filing RTIs, compiling rank-vs-marks spreadsheets, and protesting peacefully (\texttt{[06:15]}, \texttt{[15:40]}) are described with warm, admiring prosody and open-handed gestures.
  \item \textbf{Bihar Police's Patna FIR (\texttt{5 May})} detailing the arrested solver gang (\texttt{[11:10]}--\texttt{[12:30]}) is presented with crisp, factual authority, pointing at the timeline graphic on screen.
  \item \textbf{NTA's administrative actions} (changing answers in the provisional key, awarding arbitrary grace marks without prior notification in the information bulletin, \texttt{[08:10]}) are delivered with a furrowed brow and slow, incredulous head shakes.
\end{itemize}

\noindent\textbf{8.} \textbf{Acoustic \& visual space granted to students' own voices: Are protesting students heard in their own voice with clear audio, steady close-ups, and empathetic listening nods, or are they reduced to noisy background chaos B-roll while studio anchors/politicians speak over them?}\par\smallskip
\begin{itemize}
  \item \textbf{Notes:}
Mohak grants \textbf{maximum narrative primacy to students' voices}, opening the entire documentary (\texttt{[00:00]}--\texttt{[01:40]}) with direct video testimonies of NEET aspirants before he even introduces his own thesis. The student clips are played cleanly at full volume with subtitles. When transitioning from a crying student's byte to his studio camera (\texttt{[01:25]}), Mohak holds a \texttt{(1.5)} second silent pause, his face unsmiling and eyes slightly narrowed in empathy \texttt{\{\{gesture: 1.5s silent gaze + soft nod\}\}}, ensuring the student's testimony sets the emotional anchor for the entire 20-minute analysis.\par\smallskip
\end{itemize}

\noindent\textbf{9.} \textbf{Non-verbal shift between the NEET paper-leak issue and the street protest: Does the speaker's non-verbal demeanour shift from sympathetic concern (sombre tone, sighs, solidarity gestures) when discussing the exam leak itself to combative, authoritarian, or dismissive delivery once students march toward Parliament?}\par\smallskip
\begin{itemize}
  \item \textbf{Notes:}
\textbf{No hostile shift occurs.} Instead, Mohak's non-verbal delivery transitions from \textbf{intimate personal empathy} (\texttt{[00:00]}--\texttt{[02:30]}), to \textbf{rigorous data-driven investigation} (\texttt{[03:00]}--\texttt{[14:30]}), to \textbf{solemn civic warning} (\texttt{[15:30]}--\texttt{[20:12]}). When discussing students protesting on the streets and approaching the Supreme Court (\texttt{[16:10]}), his tone remains supportive and dignified, framing their protest as a desperate attempt to salvage institutional trust (\texttt{"ye sirf ek exam ki baat nahin hai (.) ye desh ke youth ke \_bharose\_ ki baat hai$\downarrow$"}, \texttt{[18:45]}).\par\smallskip
\end{itemize}

\subsubsection*{Apply the five non-verbal features (Event 2)}

\begin{longtable}{|p{0.31\textwidth}|p{0.31\textwidth}|p{0.31\textwidth}|}
\hline
Non-verbal feature & Timestamp + spoken anchor words + non-verbal notation & What it does in context (stance, framing, implication) \\ \hline \endhead
\textbf{1. Pauses, silences \& stops} & \texttt{[01:25]} \texttt{[(1.5) silent pause after student testimony]} / \texttt{[09:08]} \textit{"NTA Director ne kaha--- (1.0) 'integrity is not compromised'$\downarrow$"} & The \texttt{(1.5)} opening silence honours student grief; the \texttt{(1.0)} pause before quoting the NTA Director isolates the official claim so its contradiction with the data is glaring. \\ \hline
\textbf{2. Trailing sentences (\texttt{...})} & \texttt{[18:05]} \textit{"Jab ek \texttt{18} saal ka bachcha dekhta hai ki \texttt{2} saal ki mehnat ke baad bhi system aisa hai... toh uska bharosa..."} \texttt{\{\{gesture: slow head shake + palms dropping to lap\}\}} & Trails off on the erosion of youth trust in India's meritocracy, leaving the long-term societal consequence hanging in the air with quiet gravity. \\ \hline
\textbf{3. Intonation, pitch \& volume} & \texttt{[05:12]} \textit{"\texttt{620} marks par pichhle saal rank thi \texttt{15,000}---is saal (.) \_SAATH HAZAAR\_ (\texttt{60,000})$\uparrow$!"} / \texttt{[01:32]} \textit{"$^\circ$ek middle-class parivar apni poori jama-poonji laga deta hai$^\circ$"} & Uses a restrained baseline so that the rare pitch peak \texttt{$\uparrow$} on \texttt{SAATH HAZAAR} and the hushed drop \texttt{$^\circ$...$^\circ$} on family sacrifice carry heightened emotional and analytical contrast. \\ \hline
\textbf{4. Coughs, sighs \& vocalisations} & \texttt{[09:10]} \texttt{[[sigh]]} \textit{"Yeh press conference dekh kar answers kam mile, sawaal zyada..."} / \texttt{[09:42]} \texttt{[[chuckle]]} (on NTA appointing its own oversight panel) & The weary \texttt{[sigh]} performs disappointed civic expectation rather than partisan rage; the brief dry \texttt{[chuckle]} highlights bureaucratic absurdity. \\ \hline
\textbf{5. Hand gestures, face \& posture} & \texttt{[03:38]} \texttt{\{\{gesture: counting 3 hypotheses on fingers\}\}} / \texttt{[05:15]} \texttt{\{\{gesture: hand sweeping downward along rank-inflation graph\}\}} / \texttt{[09:10]} \texttt{\{\{gesture: slow head shake + pursed lips\}\}} & Finger-counting organises the inquiry like a research seminar; graph-tracing grounds claims in visual statistics; the slow head shake signals measured institutional critique. \\ \hline
\end{longtable}

\subsubsection*{Short summary (Event 2 Non-Verbal Synthesis)}

\begin{itemize}
  \item \textbf{What overall interpretation or emotional stance does the non-verbal delivery encourage?}
The non-verbal delivery encourages \textbf{thoughtful, data-literate empathy and deep institutional concern}. Rather than whipping up frantic anger, Mohak's calm cadence, long empathetic silences \texttt{(1.5)}, and methodical chart gestures invite the viewer to understand \textit{how} the exam system broke and \textit{why} students have every right to demand structural reform.\par\smallskip
  \item \textbf{Whose voice or physical presence is acoustically/visually foregrounded, undermined, or silenced?}
\textbf{Individual NEET aspirants (Ayushi, Anjali) and empirical score distributions} are foregrounded right from \texttt{[00:00]}, given clean audio and respectful pauses. \textbf{NTA Director General Subodh Kumar Singh's press conference assurances} are quietly undermined---not by shouting over him, but by playing his exact clip and following it with a \texttt{(1.2)} weary \texttt{[sigh]} and a devastating statistical rebuttal (\texttt{[07:18]}, \texttt{[09:10]}).\par\smallskip
  \item \textbf{What evidence supports your non-verbal reading? What complicates it (e.g., speaker baseline, format)?}
The deliberate pacing shift from \textasciitilde{}130 WPM during data analysis (\texttt{[04:00]}--\texttt{[12:00]}) down to \textasciitilde{}105 WPM with softened volume \texttt{$^\circ$...$^\circ$} and \texttt{(1.5)} silences during student impact segments (\texttt{[01:15]}, \texttt{[18:05]}) demonstrates how prosody frames the human stakes. What complicates a purely event-specific reading is \textbf{speaker baseline}: Mohak Mangal's low-decibel, seminar-style delivery and open-palm gestures are consistent across all \textit{Soch} videos, representing a deliberate counter-brand to high-decibel TV news.\par\smallskip
  \item \textbf{How do the non-verbal cues interact with the spoken words? (Do they reinforce, dramatise, soften/hedge, invert meaning via sarcasm/irony, or leave an accusation unspoken via trailing off?)}
His non-verbal cues \textbf{reinforce and visualise} complex statistical arguments (hand trajectories matching rank-inflation curves), \textbf{soften and humanise} policy analysis through quiet vocal drops \texttt{$^\circ$...$^\circ$}, and \textbf{gently invert bureaucratic jargon} via a weary \texttt{[sigh]} and dry \texttt{[chuckle]} when the NTA claims "complete transparency" (\texttt{[09:10]}).\par\smallskip
  \item \textbf{What context, speaker baseline, or format factors support or complicate this interpretation? (e.g., Is this the speaker's normal style across all videos, or a marked shift on this specific actor/claim?)}
Mohak's decision to skip Event 1 (Ram Mandir land controversy) while devoting a 20-minute data documentary to Event 2 (NEET crisis) reflects his channel's sociological baseline. Within Event 2, his restraint makes his rare non-verbal evaluations---such as the slow head shake and \texttt{[sigh]} after the NTA press conference (\texttt{[09:10]})---far more persuasive to undecided viewers than continuous shouting.\par\smallskip
  \item \textbf{How does another outlet perform the same moment, actor, or claim non-verbally?}
Compare Mohak Mangal (\texttt{[09:10]}) with \textbf{Republic TV} (\texttt{Event 2}) and \textbf{Dhruv Rathee} (\texttt{Event 2}): when addressing NTA failures, Arnab Goswami shouts at maximum volume (\texttt{WORD}) and thumps his desk, while Dhruv Rathee points an accusatory index finger at the lens with sharp prosecutorial stress. Mohak Mangal achieves an equally damning indictment of the NTA using the exact opposite non-verbal strategy: a \texttt{(1.2)} silence, a quiet \texttt{[sigh]}, and a downward-sweeping hand gesture across a scatter plot of rank inflation.\par\smallskip
  \item \textbf{My interpretation, with reasons:}
Mohak Mangal exemplifies \textbf{low-arousal, high-credibility multimodal persuasion}. By refusing to use alarmist prosody or aggressive finger-pointing, and by anchoring his video in unmuted student testimonies (\texttt{[00:00]}--\texttt{[01:40]}) and visual data graphs (\texttt{[04:10]}--\texttt{[07:45]}), his non-verbal performance depoliticises the student protest---framing the NEET controversy not as a BJP-vs-Opposition fight, but as a devastating betrayal of Indian middle-class meritocracy.\par\smallskip
\end{itemize}

\end{document}
